% Clase 13 --- La pila real y por qué "ideal" no es una propiedad de la pila
% (§3.3 y el cierre de §4.2).
% El docente descompuso la pila real en una pila ideal más una resistencia r en
% serie, y de ahí sacó la conclusión que sorprende: que una batería se comporte
% como ideal depende de A QUÉ SE LA CONECTA, no de la batería.
\begin{center}
\begin{tikzpicture}[scale=0.95]

  % =============== (a) el modelo de la pila real ===============
  \begin{scope}
    \begin{scope}[line width=0.9pt, color=figblue]
      \draw (0,0) to[battery1, l=$\varepsilon$] (0,2.2);
      \draw (0,2.2) to[R, l=$r$] (1.9,2.2);
      \draw (1.9,2.2) -- (3.6,2.2);
      \draw (3.6,2.2) to[R, l_=$R$] (3.6,0);
      \draw (3.6,0) -- (0,0);
    \end{scope}

    % la caja que es "la pila real"
    \draw[figred, dashed, line width=0.9pt] (-0.75,-0.5) rectangle (2.5,2.85);
    \node[figlbl, figred, anchor=south, align=center] at (0.85,2.9)
      {la pila real};

    \draw[figvecres] (2.5,2.55) -- (3.4,2.55);
    \node[figlbl, figred, anchor=south] at (2.95,2.6) {$I$};

    \node[figlbl, anchor=north, align=center] at (1.6,-1.55)
      {(a) pila ideal $+$ resistencia interna};
  \end{scope}

  % =============== (b) el voltaje en bornes cae con la corriente ===============
  \begin{axis}[
      fisfig,
      width=7.6cm, height=5.4cm,
      grid=none,
      axis lines=left,
      xmin=0, xmax=1.15, ymin=0, ymax=1.15,
      xlabel={$I$}, ylabel={$\Delta V$ entre bornes},
      xtick=\empty, ytick={1}, yticklabels={$\varepsilon$},
      axis line style={figgray},
      at={(7.6cm,-0.2cm)},
    ]
    \addplot[figblue, line width=1pt, domain=0:1.05] {1 - 0.75*x};
    \node[figlbl, figblue, anchor=west] at (axis cs:0.42,0.78)
      {$\Delta V = \varepsilon - r\,I$};
    \node[figlbl, figred, anchor=west, align=left] at (axis cs:0.03,0.36)
      {$I$ chica ($R \gg r$):\\[-0.2em] se porta como ideal};
    \draw[figaux] (axis cs:0,1) -- (axis cs:0.18,1);
  \end{axis}
  \node[figlbl, anchor=north, align=center] at (11.4,-1.55)
    {(b) cuánto se aparta de la ideal};

\end{tikzpicture}\\[0.5em]
{\small\itshape La fuerza electromotriz $\varepsilon$ ---que no es una fuerza
sino una energía por unidad de carga, y el nombre es sólo histórico--- es una
propiedad intrínseca de la pila: $1{,}5$ V en las comunes. Lo que \emph{no} es
intrínseco es comportarse como ideal. Toda pila disipa algo por efecto Joule en
su propio interior, y eso se modela como una resistencia $r$ en serie, de modo
que el voltaje que realmente aparece entre los bornes, $\varepsilon - rI$, baja
a medida que se le pide más corriente. Conectada a un $R$ mucho mayor que $r$,
la corriente es chica, la caída interna es despreciable y la pila parece ideal;
en cortocircuito, con $R \to 0$, la única resistencia que queda es $r$, la
corriente satura en $\varepsilon/r$ y la aproximación se rompe por completo. La
misma pila, entonces, es ideal o no según a qué se la conecte.}
\end{center}
